\documentclass[12pt,a4paper]{article}
\usepackage[utf8]{inputenc}
\usepackage{amsmath,amssymb,amsthm}
\usepackage{algorithm}
\usepackage{algpseudocode}
\usepackage{geometry}
\geometry{margin=1in}

\newtheorem{theorem}{Theorem}
\newtheorem{lemma}[theorem]{Lemma}
\newtheorem{corollary}[theorem]{Corollary}
\newtheorem{proposition}[theorem]{Proposition}

\title{Problem 70: Totient Permutation}
\author{Project Euler Solution}
\date{}

\begin{document}
\maketitle

\section{Problem Statement}

Find the value of $n$, $1 < n < 10^7$, for which $\varphi(n)$ is a permutation of $n$ and the ratio $n/\varphi(n)$ produces a minimum.

\section{Mathematical Analysis}

\begin{theorem}[Ratio Formula]
For $n = p_1^{a_1} \cdots p_k^{a_k}$:
\[
\frac{n}{\varphi(n)} = \prod_{i=1}^{k} \frac{p_i}{p_i - 1}.
\]
\end{theorem}

\begin{proof}
Immediate from $\varphi(n) = n \prod_{p \mid n}(1 - 1/p)$.
\end{proof}

\begin{theorem}[Minimization Principles]
To minimize $n/\varphi(n)$:
\begin{enumerate}
    \item Use as few distinct prime factors as possible.
    \item Use primes as large as possible.
\end{enumerate}
\end{theorem}

\begin{proof}
(1) Each distinct prime $p$ contributes a factor $p/(p-1) > 1$, so fewer factors yield a smaller product. (2) Since $p/(p-1) = 1 + 1/(p-1)$ is strictly decreasing in $p$, larger primes contribute factors closer to~1.
\end{proof}

\begin{theorem}[Exclusion of Primes]
If $p$ is prime, then $\varphi(p) = p - 1$ is never a digit permutation of $p$.
\end{theorem}

\begin{proof}
Two integers are digit permutations only if they are congruent modulo~9 (since digit sums are preserved). But $p - (p-1) = 1 \not\equiv 0 \pmod{9}$, so $p \not\equiv p-1 \pmod{9}$.
\end{proof}

\begin{corollary}[Semiprimes Are Optimal]
The candidates with fewest prime factors are semiprimes $n = pq$ ($p, q$ distinct primes). For these:
\[
\varphi(pq) = (p-1)(q-1), \quad \frac{pq}{\varphi(pq)} = \frac{pq}{(p-1)(q-1)}.
\]
This ratio is minimized when $p$ and $q$ are large and close to $\sqrt{10^7} \approx 3162$.
\end{corollary}

\begin{proof}
By Theorem~3, primes are excluded. If $n$ has $\geq 3$ distinct prime factors, then $n/\varphi(n) \geq (2/1)(3/2)(5/4) = 3.75$, which far exceeds the semiprime optimum near~1. Hence semiprimes suffice.
\end{proof}

\begin{proposition}[Search Sufficiency]
It suffices to search all pairs $(p,q)$ with $p \leq q$ prime and $pq < 10^7$.
\end{proposition}

\begin{proof}
By Theorem~3, primes are excluded. By the Corollary, numbers with $\geq 3$ prime factors are dominated. Hence only semiprimes (and prime squares $p^2$, which are included when $p = q$) need be checked.
\end{proof}

\section{Algorithm}

\begin{algorithm}[H]
\caption{Find $n < 10^7$ minimizing $n/\varphi(n)$ with $\varphi(n)$ a digit permutation of $n$}
\begin{algorithmic}[1]
\State Sieve primes up to $10^7$
\State $r_{\min} \gets \infty$, $n_{\text{best}} \gets 0$
\For{each prime $p$ with $p^2 < 10^7$}
    \For{each prime $q \ge p$ with $pq < 10^7$}
        \State $\phi \gets (p-1)(q-1)$
        \If{$\text{sorted\_digits}(pq) = \text{sorted\_digits}(\phi)$ \textbf{and} $pq / \phi < r_{\min}$}
            \State $r_{\min} \gets pq / \phi$, $n_{\text{best}} \gets pq$
        \EndIf
    \EndFor
\EndFor
\State \Return $n_{\text{best}}$
\end{algorithmic}
\end{algorithm}

\section{Complexity}

\begin{itemize}
    \item \textbf{Time}: $O(N \log \log N)$ for the sieve, plus $O(\pi(\sqrt{N})^2) = O(N/\log^2 N)$ for the search.
    \item \textbf{Space}: $O(N)$ for the sieve.
\end{itemize}

\section{Result}

The answer is $\boxed{8319823}$.

\begin{itemize}
    \item $n = 8319823 = 2339 \times 3557$
    \item $\varphi(8319823) = 2338 \times 3556 = 8313928$
    \item Sorted digits of $8319823$: $\{1, 2, 3, 3, 8, 8, 9\}$
    \item Sorted digits of $8313928$: $\{1, 2, 3, 3, 8, 8, 9\}$ \checkmark
    \item $n / \varphi(n) \approx 1.000709$
\end{itemize}

\section{Answer}

\[
\boxed{8319823}
\]

\end{document}
